% GENERATED FILE. Edit canonical lessons, then run npm run generate:books.
% canonical-source-hash: fnv1a64:3f46efe7c07601be
% canonical-lessons: GE-C199-gras, GE-C199-frucht, GE-C199-karotte, GE-C199-haehnchen, GE-C199-torte

\chapter{Grass, A fruit, A carrot, A chicken, A gateau}
\label{ch:ge-grass-a-fruit-a-carrot-a-chicken-a-gateau}
\hlchaptermodality{199}

\begin{chapteropening}
\textbf{By the end of this chapter:} \emph{I can say grass, a fruit, a carrot, a chicken and a gateau.}

Complete the last lesson of chapter 199: i can say grass, a fruit, a carrot, a chicken and a gateau.
\end{chapteropening}

\section[das Gras]{das Gras — grass}

\label{lesson:GE-C199-gras}



\begin{quote}
Before the new one: say the German for a butterfly, then the German for a plant.
\end{quote}

\subsection*{You'll want to know: das Gras}
\textbf{das Gras} — \textquotedblleft{}grass\textquotedblright{}.

\begin{grammarlens}[title={What you've built}]
One more word for this chapter.
\end{grammarlens}

\subsection*{Your turn}
\begin{itemize}
\raggedright
  \item \emph{Say it:} \emph{das Gras}
  \item \emph{Say it:} \emph{das Gras}, once more
  \item \emph{Say it:} say \emph{das Gras}
  \item \emph{Recall:} say the German for a butterfly, then the German for a plant, then say \emph{das Gras} again
\end{itemize}

\begin{tcolorbox}[breakable,colback=teal!4,colframe=teal!35!black,title={Before you move on}]
What does das Gras mean? (\textquotedblleft{}Grass\textquotedblright{}.) Say it once more.
\end{tcolorbox}
\section[die Frucht]{die Frucht — a fruit}

\label{lesson:GE-C199-frucht}



\begin{quote}
Before the new one: say the German for a plant, then the German for grass.
\end{quote}

\subsection*{You'll want to know: die Frucht}
\textbf{die Frucht} — \textquotedblleft{}a fruit\textquotedblright{}.

\begin{grammarlens}[title={What you've built}]
One more word for this chapter.
\end{grammarlens}

\subsection*{Your turn}
\begin{itemize}
\raggedright
  \item \emph{Say it:} \emph{die Frucht}
  \item \emph{Say it:} \emph{die Frucht}, once more
  \item \emph{Say it:} say \emph{die Frucht}
  \item \emph{Recall:} say the German for a plant, then the German for grass, then say \emph{die Frucht} again
\end{itemize}

\begin{tcolorbox}[breakable,colback=teal!4,colframe=teal!35!black,title={Before you move on}]
What does die Frucht mean? (\textquotedblleft{}A fruit\textquotedblright{}.) Say it once more.
\end{tcolorbox}
\section[die Karotte]{die Karotte — a carrot}

\label{lesson:GE-C199-karotte}



\begin{quote}
Before the new one: say the German for grass, then the German for a fruit.
\end{quote}

\subsection*{You'll want to know: die Karotte}
\textbf{die Karotte} — \textquotedblleft{}a carrot\textquotedblright{}.

\begin{grammarlens}[title={What you've built}]
One more word for this chapter.
\end{grammarlens}

\subsection*{Your turn}
\begin{itemize}
\raggedright
  \item \emph{Say it:} \emph{die Karotte}
  \item \emph{Say it:} \emph{die Karotte}, once more
  \item \emph{Say it:} say \emph{die Karotte}
  \item \emph{Recall:} say the German for grass, then the German for a fruit, then say \emph{die Karotte} again
\end{itemize}

\begin{tcolorbox}[breakable,colback=teal!4,colframe=teal!35!black,title={Before you move on}]
What does die Karotte mean? (\textquotedblleft{}A carrot\textquotedblright{}.) Say it once more.
\end{tcolorbox}
\section[das Hähnchen]{das Hähnchen — a chicken (to eat)}

\label{lesson:GE-C199-haehnchen}



\begin{quote}
Before the new one: say the German for a fruit, then the German for a carrot.
\end{quote}

\subsection*{You'll want to know: das Hähnchen}
\textbf{das Hähnchen} — \textquotedblleft{}a chicken (to eat)\textquotedblright{}.

\begin{grammarlens}[title={What you've built}]
One more word for this chapter.
\end{grammarlens}

\subsection*{Your turn}
\begin{itemize}
\raggedright
  \item \emph{Say it:} \emph{das Hähnchen}
  \item \emph{Say it:} \emph{das Hähnchen}, once more
  \item \emph{Say it:} say \emph{das Hähnchen}
  \item \emph{Recall:} say the German for a fruit, then the German for a carrot, then say \emph{das Hähnchen} again
\end{itemize}

\begin{tcolorbox}[breakable,colback=teal!4,colframe=teal!35!black,title={Before you move on}]
What does das Hähnchen mean? (\textquotedblleft{}A chicken (to eat)\textquotedblright{}.) Say it once more.
\end{tcolorbox}
\section[die Torte]{die Torte — a gateau, a layer cake}

\label{lesson:GE-C199-torte}



\begin{quote}
Before the new one: say the German for a carrot, then the German for a chicken.
\end{quote}

\subsection*{You'll want to know: die Torte}
\textbf{die Torte} — \textquotedblleft{}a gateau, a layer cake\textquotedblright{}.

\begin{grammarlens}[title={What you've built}]
One more word for this chapter.
\end{grammarlens}

\subsection*{Your turn}
\begin{itemize}
\raggedright
  \item \emph{Say it:} \emph{die Torte}
  \item \emph{Say it:} \emph{die Torte}, once more
  \item \emph{Say it:} say \emph{die Torte}
  \item \emph{Recall:} say the German for a carrot, then the German for a chicken, then say \emph{die Torte} again
\end{itemize}

\begin{tcolorbox}[breakable,colback=teal!4,colframe=teal!35!black,title={Before you move on}]
What does die Torte mean? (\textquotedblleft{}A gateau, a layer cake\textquotedblright{}.) Say it once more.
\end{tcolorbox}
